% =====================================================
% part5_inner.tex ― 第5部:内積空間
% =====================================================
\part{内積空間}

% =====================================================
\chapter{内積とノルム}\label{chap:inner}
% =====================================================

\begin{goal}
\begin{itemize}
\item $\R^n$ の標準内積とノルムを計算でき、コーシー・シュワルツの不等式と
三角不等式を\textbf{証明つきで}理解する。
\item 内積を使って $n$ 次元でも「角度」と「直交」が定義できることを理解し、
直交条件・距離・単位ベクトルの計算を実行できるようになる。
\item 内積の\textbf{公理}(対称・双線形・正値)を知り、関数の内積のような
$\R^n$ 以外の内積空間でも同じ定理が使えることを理解する。
\end{itemize}
\end{goal}

\section{動機 ― 長さと角度を舞台に呼び戻す}

第IV部までの理論には、実は「長さ」と「角度」が登場しなかった。
和とスカラー倍だけの世界(線形構造)で、どこまで行けるかを見てきたのである。
基底も次元も固有値も、長さを測らずに定義できた。
本部では幾何構造――内積――を導入する。
舞台に定規と分度器が持ち込まれ、「直交」という強力な概念が使えるようになる。

出発点は高校の内積である。平面ベクトルについて
\[
\vect{a} \cdot \vect{b} = |\vect{a}|\,|\vect{b}| \cos\theta
= a_1 b_1 + a_2 b_2
\]
と、「長さと角度による定義」と「成分による計算式」の二本立てで学んだ。
$n$ 次元ではどちらを定義に採用すべきか。$4$ 次元以上では
「角度」を先に測る手段がないから、\textbf{成分の式を定義に採用し、
そこから逆に角度を定義する}のが正しい順序である。
この逆転がうまくいくこと(角度が矛盾なく定義できること)を保証するのが、
本章の主役コーシー・シュワルツの不等式である。

\section{標準内積とノルム}

\begin{teigi}[標準内積とノルム]\label{def:inner-std}
$\vect{a}, \vect{b} \in \R^n$ に対し、
\[
\langle \vect{a}, \vect{b} \rangle
= a_1 b_1 + a_2 b_2 + \cdots + a_n b_n = {}^t\vect{a}\,\vect{b}
\]
を\textbf{標準内積}という。また
\[
\| \vect{a} \| = \sqrt{\langle \vect{a}, \vect{a} \rangle}
= \sqrt{a_1^2 + \cdots + a_n^2}
\]
を $\vect{a}$ の\textbf{ノルム}(長さ)という。
\end{teigi}

$n = 2, 3$ では高校の内積・大きさと完全に一致する。
${}^t\vect{a}\,\vect{b}$ という行列の積の形($1 \times n$ 行列と
$n \times 1$ 行列の積は $1 \times 1$、すなわち数)で書けることは、
転置と積の計算法則(定理 \ref{thm:transpose})を内積に直結させる
便利な表記であり、本部で繰り返し使う。

\begin{teiri}[内積の基本性質]\label{thm:inner-props}
任意の $\vect{a}, \vect{b}, \vect{c} \in \R^n$ と $k \in \R$ に対して:
\begin{enumerate}
\item \textbf{対称性}:$\langle \vect{a}, \vect{b} \rangle = \langle \vect{b}, \vect{a} \rangle$。
\item \textbf{線形性}:$\langle \vect{a} + \vect{b}, \vect{c} \rangle
= \langle \vect{a}, \vect{c} \rangle + \langle \vect{b}, \vect{c} \rangle$、
$\langle k\vect{a}, \vect{b} \rangle = k\langle \vect{a}, \vect{b} \rangle$
(対称性より第 $2$ 変数についても同様。あわせて\textbf{双線形}という)。
\item \textbf{正値性}:$\langle \vect{a}, \vect{a} \rangle \ge 0$ であり、
等号は $\vect{a} = \vect{0}$ のときに限る。
\end{enumerate}
\end{teiri}

\noindent\textbf{(証明)}\quad
すべて成分の計算に帰着する。
(1) $a_i b_i = b_i a_i$(実数の積の交換律)を足し合わせればよい。
(2) 第 $i$ 項が $(a_i + b_i)c_i = a_ic_i + b_ic_i$、$(ka_i)b_i = k(a_ib_i)$
となること(実数の分配律・結合律)から従う。
(3) $\langle \vect{a}, \vect{a} \rangle = a_1^2 + \cdots + a_n^2$ は
$0$ 以上の実数の和だから $0$ 以上で、和が $0$ になるのは
各 $a_i^2 = 0$、すなわちすべての成分が $0$ のときに限る。 $\blacksquare$

ノルムについては、定義からただちに
$\|k\vect{a}\| = |k|\,\|\vect{a}\|$ と、
$\|\vect{a}\| \ge 0$(等号は $\vect{a} = \vect{0}$ のみ)が分かる。
ノルム $1$ のベクトルを\textbf{単位ベクトル}という。
$\vect{a} \neq \vect{0}$ を自分のノルムで割った
$\dfrac{1}{\|\vect{a}\|}\vect{a}$ は、$\vect{a}$ と同じ向きの単位ベクトルである
(この操作を\textbf{正規化}という)。また
\[
d(\vect{a}, \vect{b}) = \|\vect{a} - \vect{b}\|
\]
を $2$ 点(の位置ベクトル)$\vect{a}, \vect{b}$ の\textbf{距離}という。
$n = 2, 3$ では座標幾何の $2$ 点間距離の公式そのものである。

\begin{rei}[計算例]\label{rei:inner-firstcalc}
$\vect{a} = \mat{1 \\ 2 \\ 2}$ について
$\|\vect{a}\| = \sqrt{1 + 4 + 4} = 3$。正規化すると
$\dfrac{1}{3}\mat{1 \\ 2 \\ 2}$(ノルム $1$ を確かめよ)。
$\vect{b} = \mat{2 \\ -2 \\ 1}$ との内積は
$\langle \vect{a}, \vect{b} \rangle = 2 - 4 + 2 = 0$、距離は
$d(\vect{a}, \vect{b}) = \|(-1, 4, 1)\| = \sqrt{18} = 3\sqrt{2}$ である。
\end{rei}

\section{コーシー・シュワルツの不等式と三角不等式}

\begin{teiri}[コーシー・シュワルツの不等式]\label{thm:inner-cs}
任意の $\vect{a}, \vect{b} \in \R^n$ に対して
\[
|\langle \vect{a}, \vect{b} \rangle| \le \|\vect{a}\|\,\|\vect{b}\|.
\]
等号が成り立つのは、$\vect{a}, \vect{b}$ の一方が他方のスカラー倍
(線形従属)のとき、かつそのときに限る。
\end{teiri}

\noindent\textbf{(証明)}\quad
高校で習った判別式の技がそのまま使える。
$\vect{a} = \vect{0}$ なら両辺とも $0$ で成立する(等号の場合も
$\vect{0} = 0\vect{b}$ で線形従属)。$\vect{a} \neq \vect{0}$ とし、
任意の実数 $t$ に対して
\[
f(t) = \| t\vect{a} + \vect{b} \|^2
= t^2 \|\vect{a}\|^2 + 2t\langle \vect{a}, \vect{b} \rangle + \|\vect{b}\|^2
\]
を考える(展開には定理 \ref{thm:inner-props} の双線形性を使った)。
ノルムの $2$ 乗だから $f(t) \ge 0$ が\textbf{すべての} $t$ で成り立つ。
$t^2$ の係数 $\|\vect{a}\|^2 > 0$ の $2$ 次式が常に非負だから判別式は
$0$ 以下、すなわち
\[
\langle \vect{a}, \vect{b} \rangle^2 - \|\vect{a}\|^2 \|\vect{b}\|^2 \le 0
\]
であり、平方根をとれば不等式を得る。
等号が成り立つのは判別式が $0$、すなわちある $t_0$ で $f(t_0) = 0$ と
なるときである。$f(t_0) = \|t_0\vect{a} + \vect{b}\|^2 = 0$ は正値性より
$\vect{b} = -t_0\vect{a}$ を意味し、これは線形従属にほかならない。
逆に $\vect{b} = k\vect{a}$ なら両辺とも $|k|\,\|\vect{a}\|^2$ で等号が成り立つ。
$\blacksquare$

この不等式のおかげで、$n$ 次元でも
\[
\cos\theta = \frac{\langle \vect{a}, \vect{b} \rangle}{\|\vect{a}\|\,\|\vect{b}\|}
\qquad (\vect{a}, \vect{b} \neq \vect{0})
\]
の右辺が必ず $-1$ 以上 $1$ 以下に収まり、\textbf{角度 $\theta$
($0 \le \theta \le 180^\circ$)が矛盾なく定義できる}。
測れないはずの $4$ 次元以上の「角度」が、内積を経由して意味をもつ。

\begin{teigi}[直交]\label{def:inner-orth}
$\langle \vect{a}, \vect{b} \rangle = 0$ のとき、$\vect{a}$ と $\vect{b}$ は
\textbf{直交}するといい、$\vect{a} \perp \vect{b}$ と書く。
零ベクトルは(角度は定義されないが)すべてのベクトルと直交すると約束する。
\end{teigi}

\begin{teiri}[三角不等式]\label{thm:inner-triangle}
任意の $\vect{a}, \vect{b} \in \R^n$ に対して
\[
\|\vect{a} + \vect{b}\| \le \|\vect{a}\| + \|\vect{b}\|.
\]
\end{teiri}

\noindent\textbf{(証明)}\quad
両辺とも $0$ 以上だから、$2$ 乗して比べればよい。
コーシー・シュワルツの不等式(定理 \ref{thm:inner-cs})より
$\langle \vect{a}, \vect{b} \rangle \le \|\vect{a}\|\,\|\vect{b}\|$ だから
\[
\|\vect{a} + \vect{b}\|^2
= \|\vect{a}\|^2 + 2\langle \vect{a}, \vect{b} \rangle + \|\vect{b}\|^2
\le \|\vect{a}\|^2 + 2\|\vect{a}\|\,\|\vect{b}\| + \|\vect{b}\|^2
= \bigl(\|\vect{a}\| + \|\vect{b}\|\bigr)^2. \qquad\blacksquare
\]

「寄り道すると遠くなる」という平面図形の常識(三角形の $2$ 辺の和は
他の $1$ 辺より長い)が、$n$ 次元でも成り立つわけである。
コーシー・シュワルツの不等式が幾何の常識を高次元へ運ぶ橋になっている。

\begin{teiri}[ピタゴラスの定理]\label{thm:inner-pyth}
$\vect{a} \perp \vect{b}$ ならば
$\|\vect{a} + \vect{b}\|^2 = \|\vect{a}\|^2 + \|\vect{b}\|^2$。
\end{teiri}

\noindent\textbf{(証明)}\quad
$\|\vect{a} + \vect{b}\|^2 = \|\vect{a}\|^2 + 2\langle \vect{a}, \vect{b} \rangle
+ \|\vect{b}\|^2$ で、直交より中央の項が $0$。 $\blacksquare$

三平方の定理が $n$ 次元で(しかも逆向きも。問 \ref{q:inner-ex12})
成り立つ。次章で正射影が「最も近い点」であることを示すとき、
この定理が決定打になる。

\begin{itembox}[l]{深掘り:相関係数の正体は $\cos\theta$ である}
統計学で学ぶ相関係数 $r$ は、$2$ 種類のデータ
$x_1, \dots, x_n$ と $y_1, \dots, y_n$ の関係の強さを測る量で、
「必ず $-1 \le r \le 1$ になる」と教わる。なぜか。
各データから平均を引いた偏差ベクトル
$\vect{x} = (x_1 - \bar{x}, \dots, x_n - \bar{x})$,
$\vect{y} = (y_1 - \bar{y}, \dots, y_n - \bar{y})$ を作ると、
相関係数の定義式は
\[
r = \frac{\langle \vect{x}, \vect{y} \rangle}{\|\vect{x}\|\,\|\vect{y}\|}
\]
と、\textbf{$n$ 次元空間における $2$ 本のデータベクトルのなす角の余弦}
そのものである。$|r| \le 1$ はコーシー・シュワルツの不等式であり、
$r = \pm 1$(等号成立)はデータベクトルが平行、
すなわち完全な線形関係があることを意味する。
「無相関」とはデータベクトルが\textbf{直交}していることだ
――統計の概念が、幾何の言葉に翻訳される瞬間である。
\end{itembox}

\section{抽象的な内積空間}

第III部でベクトルを公理化したのと同様、内積も公理化できる。
定理 \ref{thm:inner-props} で\textbf{証明した}性質を、こんどは
\textbf{要求する公理}に格上げするのである。

\begin{teigi}[内積空間]\label{def:inner-axioms}
ベクトル空間 $V$ の任意の $2$ 元 $\vect{u}, \vect{v}$ に実数
$\langle \vect{u}, \vect{v} \rangle$ を対応させる演算が
\begin{enumerate}
\item 対称性:$\langle \vect{u}, \vect{v} \rangle = \langle \vect{v}, \vect{u} \rangle$、
\item 双線形性:$\langle \vect{u} + \vect{u}', \vect{v} \rangle
= \langle \vect{u}, \vect{v} \rangle + \langle \vect{u}', \vect{v} \rangle$、
$\langle k\vect{u}, \vect{v} \rangle = k\langle \vect{u}, \vect{v} \rangle$
(第 $2$ 変数も同様)、
\item 正値性:$\langle \vect{u}, \vect{u} \rangle \ge 0$、等号は
$\vect{u} = \vect{0}$ のみ、
\end{enumerate}
を満たすとき、この演算を\textbf{内積}、組 $(V, \langle\ ,\ \rangle)$ を
\textbf{内積空間}という。ノルム・距離・角度・直交は、標準内積のときと
同じ式で定義する。
\end{teigi}

\begin{chui}
コーシー・シュワルツの不等式(定理 \ref{thm:inner-cs})と
三角不等式(定理 \ref{thm:inner-triangle})の証明を見返すと、
使ったのは\textbf{双線形性と正値性だけ}で、成分は一度も登場していない。
したがって両定理は、どんな内積空間でも\textbf{再証明なしに}成り立つ。
第III部で学んだ「公理からの証明は使い回せる」という節約が、
ここでも効いている。
\end{chui}

重要な例を一つ挙げる。

\begin{rei}[関数の内積]\label{rei:inner-func}
区間 $[-1, 1]$ 上の連続関数の空間で、
\[
\langle f, g \rangle = \int_{-1}^{1} f(x)\,g(x)\,dx
\]
は内積の公理を満たす。たとえば $f(x) = 1$ と $g(x) = x$ は
\[
\langle 1, x \rangle = \int_{-1}^{1} x\,dx = 0
\]
なので\textbf{直交}している。「関数どうしが直交する」という、
高校では意味をなさなかった文が、ここでは厳密な意味をもつ。
$\sin nx$ と $\cos mx$ たちの直交性を利用して任意の周期関数を
三角関数の「直交基底」で展開する理論が\textbf{フーリエ級数}であり、
音声処理から画像圧縮、量子力学まで現代技術の土台になっている。
本書では立ち入らないが、その入口がここにあることを覚えておいてほしい。
\end{rei}

\begin{mondai}\label{q:innerp}
　(1) $\vect{a} = \mat{1 \\ 1 \\ 0}$, $\vect{b} = \mat{1 \\ 0 \\ 1}$ の
なす角を求めよ。\\
　(2) 上の関数の内積に関して、$f(x) = x$ と $g(x) = x^2$ が
直交することを示せ。
\end{mondai}

\section{反例 ― ありがちな誤解を正す}

\begin{hanrei}[正値性が破れると「長さ」が壊れる]\label{hanrei:inner-indefinite}
$\R^2$ に、内積もどきの演算
\[
\langle \vect{a}, \vect{b} \rangle' = a_1 b_1 - a_2 b_2
\]
を定めてみる(第 $2$ 成分の符号を反転した)。対称性と双線形性は
成分計算でそのまま確かめられるが、\textbf{正値性が破れる}。実際
\[
\vect{a} = \mat{0 \\ 1} \ \text{に対し}\
\langle \vect{a}, \vect{a} \rangle' = 0 - 1 = -1 < 0,
\qquad
\vect{b} = \mat{1 \\ 1} \ \text{に対し}\
\langle \vect{b}, \vect{b} \rangle' = 1 - 1 = 0
\ (\vect{b} \neq \vect{0}).
\]
「長さの $2$ 乗」が負になったり、零ベクトルでないのに $0$ に
なったりするので、この演算からはノルムが定義できず、
コーシー・シュワルツの不等式も保証されない。内積の三つの公理は
どれも飾りではない。
(なお、この符号ちがいの演算そのものは無意味ではなく、
相対性理論の時空の幾何(ミンコフスキー計量)で主役を務める。
「正値性を外すと何が起こるか」を真剣に研究した理論である。)
\end{hanrei}

\begin{hanrei}[内積は「消去」できない]\label{hanrei:inner-cancel}
数の等式 $ab = ac$($a \neq 0$)なら $b = c$ と消去できるが、
内積では $\langle \vect{a}, \vect{b} \rangle = \langle \vect{a}, \vect{c} \rangle$
かつ $\vect{a} \neq \vect{0}$ でも $\vect{b} = \vect{c}$ とは限らない。
\[
\vect{a} = \mat{1 \\ 0}, \quad \vect{b} = \mat{0 \\ 1}, \quad
\vect{c} = \mat{0 \\ 2}
\]
とすると $\langle \vect{a}, \vect{b} \rangle = \langle \vect{a}, \vect{c} \rangle = 0$
だが $\vect{b} \neq \vect{c}$ である。内積は $\vect{a}$ 方向の情報しか
見ておらず、それと直交する方向の違いを検出できないからである。
ただし、\textbf{すべての} $\vect{a}$ に対して
$\langle \vect{a}, \vect{b} \rangle = \langle \vect{a}, \vect{c} \rangle$
が成り立つなら $\vect{b} = \vect{c}$ といえる(問 \ref{q:inner-ex9})。
\end{hanrei}

\section{例題}

内積・ノルム・角度の計算と、不等式の使い方を練習する。
以下は\textbf{解答つきの例題}である。手を動かして追ってほしい。

\begin{rei}[例題・基本(計算)]\label{rei:inner-angle}
$\R^4$ のベクトル $\vect{a} = (1, 1, 1, 1)$, $\vect{b} = (1, 1, 1, -1)$ の
なす角を求めよ。\\
\textbf{解.}\ $\langle \vect{a}, \vect{b} \rangle = 1 + 1 + 1 - 1 = 2$、
$\|\vect{a}\| = \|\vect{b}\| = \sqrt{4} = 2$。よって
\[
\cos\theta = \frac{2}{2 \cdot 2} = \frac{1}{2}, \qquad \theta = 60^\circ.
\]
目で見ることのできない $4$ 次元のベクトルにも、内積を経由すれば
角度が確定する。これが「成分の式を定義に採用する」ことの威力である。
\end{rei}

\begin{rei}[例題・基本(直交条件)]\label{rei:inner-orthparam}
$\vect{a} = \mat{2 \\ -1 \\ 3}$ と $\vect{b} = \mat{1 \\ t \\ 2}$ が
直交するように実数 $t$ を定めよ。\\
\textbf{解.}\ 直交条件は $\langle \vect{a}, \vect{b} \rangle = 0$ だから
\[
2 \cdot 1 + (-1) \cdot t + 3 \cdot 2 = 8 - t = 0, \qquad t = 8.
\]
検算:$t = 8$ のとき $\langle \vect{a}, \vect{b} \rangle = 2 - 8 + 6 = 0$。
「垂直」という幾何の条件が、内積 $= 0$ という $1$ 本の方程式に翻訳される。
\end{rei}

\begin{rei}[例題・基本(単位ベクトルと距離)]\label{rei:inner-normdist}
　(1) $\vect{v} = \mat{3 \\ -4}$ と同じ向きの単位ベクトルを求めよ。\\
　(2) $2$ 点 $\vect{a} = \mat{1 \\ 2}$, $\vect{b} = \mat{4 \\ 6}$ の距離を求めよ。\\
\textbf{解.}\ (1) $\|\vect{v}\| = \sqrt{9 + 16} = 5$ だから、正規化して
$\dfrac{1}{5}\mat{3 \\ -4} = \mat{3/5 \\ -4/5}$。
ノルムを確かめると $\sqrt{9/25 + 16/25} = 1$ で正しい。\\
　(2) $\vect{a} - \vect{b} = \mat{-3 \\ -4}$ より
$d(\vect{a}, \vect{b}) = \sqrt{9 + 16} = 5$。
座標平面の $2$ 点間距離の公式と同じ計算である。
\end{rei}

\begin{rei}[例題・標準(コーシー・シュワルツの応用)]\label{rei:inner-csapp}
任意の実数 $a, b, c$ に対して
\[
(a + b + c)^2 \le 3(a^2 + b^2 + c^2)
\]
が成り立つことを示し、等号成立の条件を求めよ。\\
\textbf{解.}\ $\vect{x} = (a, b, c)$, $\vect{y} = (1, 1, 1)$ にコーシー・
シュワルツの不等式(定理 \ref{thm:inner-cs})を使う。
$\langle \vect{x}, \vect{y} \rangle = a + b + c$、
$\|\vect{x}\|^2 = a^2 + b^2 + c^2$、$\|\vect{y}\|^2 = 3$ だから
\[
(a + b + c)^2 = \langle \vect{x}, \vect{y} \rangle^2
\le \|\vect{x}\|^2\,\|\vect{y}\|^2 = 3(a^2 + b^2 + c^2).
\]
等号は $\vect{x}$ と $\vect{y}$ が線形従属、すなわち
$(a, b, c) = k(1, 1, 1)$、つまり $a = b = c$ のとき(逆も明らか)。
高校で「うまい変形」を要した不等式が、ベクトルを一組選ぶだけで落ちる。
\end{rei}

\begin{rei}[例題・標準(関数の内積)]\label{rei:inner-funccalc}
区間 $[-1, 1]$ 上の内積 $\langle f, g \rangle = \int_{-1}^{1} f(x)g(x)\,dx$
について、$f(x) = 1$, $g(x) = x^2$ の内積とそれぞれのノルムを計算し、
コーシー・シュワルツの不等式が成り立っていることを確かめよ。\\
\textbf{解.}\
\[
\langle 1, x^2 \rangle = \int_{-1}^{1} x^2\,dx = \frac{2}{3}, \qquad
\|1\|^2 = \int_{-1}^{1} 1\,dx = 2, \qquad
\|x^2\|^2 = \int_{-1}^{1} x^4\,dx = \frac{2}{5}.
\]
よって $\|1\| = \sqrt{2}$, $\|x^2\| = \sqrt{2/5}$。不等式の両辺を
$2$ 乗で比べると
\[
\langle 1, x^2 \rangle^2 = \frac{4}{9} \le \|1\|^2\|x^2\|^2
= 2 \cdot \frac{2}{5} = \frac{4}{5}
\]
で、確かに成り立っている($4/9 < 4/5$)。等号でないのは、$1$ と $x^2$ が
線形従属でない(定数倍の関係にない)ことの反映である。
\end{rei}

\begin{rei}[例題・標準(証明:平行四辺形の法則)]\label{rei:inner-parallelogram}
任意の $\vect{a}, \vect{b} \in \R^n$ に対して
\[
\|\vect{a} + \vect{b}\|^2 + \|\vect{a} - \vect{b}\|^2
= 2\|\vect{a}\|^2 + 2\|\vect{b}\|^2
\]
が成り立つことを示せ(平行四辺形の $2$ 本の対角線の $2$ 乗和は、
$4$ 辺の $2$ 乗和に等しい)。\\
\textbf{解.}\ 双線形性(定理 \ref{thm:inner-props})で展開する。
\[
\|\vect{a} + \vect{b}\|^2 = \|\vect{a}\|^2 + 2\langle \vect{a}, \vect{b} \rangle
+ \|\vect{b}\|^2, \qquad
\|\vect{a} - \vect{b}\|^2 = \|\vect{a}\|^2 - 2\langle \vect{a}, \vect{b} \rangle
+ \|\vect{b}\|^2.
\]
辺々加えると内積の項が打ち消し合い、
$2\|\vect{a}\|^2 + 2\|\vect{b}\|^2$ を得る。
成分を書かずに、双線形性という「計算規則」だけで証明が済む点に注目せよ。
この等式は標準内積に限らず、どんな内積空間のノルムでも成り立つ。
\end{rei}

\begin{matome}
\begin{itemize}
\item \textbf{標準内積}は $\langle \vect{a}, \vect{b} \rangle = \sum a_ib_i
= {}^t\vect{a}\,\vect{b}$、\textbf{ノルム}は
$\|\vect{a}\| = \sqrt{\langle \vect{a}, \vect{a} \rangle}$。
対称・双線形・正値という三つの基本性質をもつ(定理 \ref{thm:inner-props})。
\item \textbf{コーシー・シュワルツの不等式}
$|\langle \vect{a}, \vect{b} \rangle| \le \|\vect{a}\|\,\|\vect{b}\|$
(等号は線形従属のとき)により、$n$ 次元でも
$\cos\theta$ を通じて\textbf{角度}が定義できる。
$\langle \vect{a}, \vect{b} \rangle = 0$ が\textbf{直交}である。
\item コーシー・シュワルツから\textbf{三角不等式}
$\|\vect{a} + \vect{b}\| \le \|\vect{a}\| + \|\vect{b}\|$ が従う。
直交すれば\textbf{ピタゴラスの定理}が使える。
\item 対称・双線形・正値を公理とする\textbf{内積空間}では、
上の定理がすべて再証明なしに成り立つ。関数の空間の
$\langle f, g \rangle = \int f g\,dx$ が代表例で、
「関数どうしの直交」が意味をもつ。
\item 正値性の破れた演算(反例 \ref{hanrei:inner-indefinite})では
長さが壊れる。また内積の値が一致しても、ベクトルは消去できない
(反例 \ref{hanrei:inner-cancel})。
\end{itemize}
\end{matome}

\section*{章末問題}

\begin{mondai}[基本]\label{q:inner-ex1}
次の各組のベクトルの内積・ノルム・なす角を求めよ。\\
　(1) $\vect{a} = \mat{2 \\ 1}$, $\vect{b} = \mat{3 \\ -1}$\\
　(2) $\vect{a} = \mat{1 \\ 2 \\ 2}$, $\vect{b} = \mat{-2 \\ 1 \\ 2}$
(なす角は $\cos\theta$ の値まででよい)
\end{mondai}

\begin{mondai}[基本]\label{q:inner-ex2}
$\vect{a} = \mat{1 \\ -2 \\ 3}$ と $\vect{b} = \mat{2 \\ t \\ 4}$ が
直交するように実数 $t$ を定めよ。
\end{mondai}

\begin{mondai}[基本]\label{q:inner-ex3}
　(1) $\vect{v} = \mat{1 \\ 2 \\ -2}$ と同じ向きの単位ベクトルを求めよ。\\
　(2) $2$ 点 $\vect{a} = \mat{3 \\ 1}$, $\vect{b} = \mat{0 \\ 5}$ の
距離を求めよ。
\end{mondai}

\begin{mondai}[基本]\label{q:inner-ex4}
区間 $[-1, 1]$ 上の内積 $\langle f, g \rangle = \int_{-1}^{1} f(x)g(x)\,dx$
について、$f(x) = 1 + x$, $g(x) = x$ の内積 $\langle f, g \rangle$ と
ノルム $\|f\|$ を求めよ。
\end{mondai}

\begin{mondai}[標準]\label{q:inner-ex5}
コーシー・シュワルツの不等式を使って、任意の実数 $a, b, c$ に対して
\[
(a + 2b + 2c)^2 \le 9(a^2 + b^2 + c^2)
\]
が成り立つことを示し、等号成立の条件を求めよ。
\end{mondai}

\begin{mondai}[標準]\label{q:inner-ex6}
$\|\vect{a}\| = 2$, $\|\vect{b}\| = 3$,
$\langle \vect{a}, \vect{b} \rangle = -4$ のとき、\\
　(1) $\|\vect{a} + \vect{b}\|$ と $\|\vect{a} - \vect{b}\|$ を求めよ
(成分は与えられていないことに注意)。\\
　(2) $\|\vect{a}\| = 2$, $\|\vect{b}\| = 3$ のまま
$\langle \vect{a}, \vect{b} \rangle = 7$ となることはあり得るか。理由をつけて答えよ。
\end{mondai}

\begin{mondai}[標準]\label{q:inner-ex7}
$\R^2$ に演算 $\langle \vect{a}, \vect{b} \rangle_w = 2a_1b_1 + 3a_2b_2$ を
定める。\\
　(1) これが内積の公理(定義 \ref{def:inner-axioms})を満たすことを確かめよ。\\
　(2) この内積での $\vect{e}_1 = \mat{1 \\ 0}$ のノルムを求めよ
(標準内積とは「長さ」の値が変わることに注意)。
\end{mondai}

\begin{mondai}[標準]\label{q:inner-ex8}
区間 $[-1, 1]$ 上の内積 $\langle f, g \rangle = \int_{-1}^{1} f(x)g(x)\,dx$
について、$x^2 - c$ が定数関数 $1$ と直交するように定数 $c$ を定めよ。
\end{mondai}

\begin{mondai}[発展]\label{q:inner-ex9}
$\R^n$ のベクトル $\vect{b}, \vect{c}$ について、\textbf{すべての}
$\vect{a} \in \R^n$ に対して $\langle \vect{a}, \vect{b} \rangle
= \langle \vect{a}, \vect{c} \rangle$ が成り立つならば
$\vect{b} = \vect{c}$ であることを示せ
(ヒント:$\vect{a} = \vect{b} - \vect{c}$ を代入する)。
\end{mondai}

\begin{mondai}[発展]\label{q:inner-ex10}
$n$ 次元立方体の対角線ベクトル $\vect{d} = (1, 1, \dots, 1) \in \R^n$ と
辺ベクトル $\vect{e}_1 = (1, 0, \dots, 0)$ のなす角を $\theta_n$ とする。\\
　(1) $n = 4$ のとき $\theta_4$ を求めよ。\\
　(2) 一般の $n$ で $\cos\theta_n$ を求め、$n$ を大きくすると
$\theta_n$ がどうなるか述べよ。
\end{mondai}

\begin{mondai}[証明]\label{q:inner-ex11}
　(1) 三角不等式(定理 \ref{thm:inner-triangle})を用いて、
\textbf{逆向き三角不等式}
$\bigl|\,\|\vect{a}\| - \|\vect{b}\|\,\bigr| \le \|\vect{a} - \vect{b}\|$
を証明せよ
(ヒント:$\vect{a} = (\vect{a} - \vect{b}) + \vect{b}$ と分解する)。\\
　(2) 等号が成り立つ例を一つ挙げよ。
\end{mondai}

\begin{mondai}[証明]\label{q:inner-ex12}
ピタゴラスの定理の\textbf{逆}、すなわち
「$\|\vect{a} + \vect{b}\|^2 = \|\vect{a}\|^2 + \|\vect{b}\|^2$ ならば
$\vect{a} \perp \vect{b}$」を証明せよ。
また、この逆が成り立つ理由を、展開式のどの項が消えることと対応するか
説明せよ。
\end{mondai}

\bigskip
{\small
\noindent\textbf{【章末問題の方針】}\quad
まず自力で取り組み、行き詰まったら下の方針を見よ。記述による完全解答は
巻末「問の解答」にある。
\begin{itemize}
\item \textbf{問 \ref{q:inner-ex1}}:内積・ノルムを定義どおり計算し、
$\cos\theta = \langle \vect{a}, \vect{b} \rangle / (\|\vect{a}\|\|\vect{b}\|)$
に代入する。
\item \textbf{問 \ref{q:inner-ex2}}:直交 $\iff$ 内積 $= 0$。$t$ の
$1$ 次方程式になる。
\item \textbf{問 \ref{q:inner-ex3}}:単位ベクトルはノルムで割る。距離は
差のノルム。
\item \textbf{問 \ref{q:inner-ex4}}:定義の積分を計算する。
奇関数の対称区間の積分は $0$ になることを使うと速い。
\item \textbf{問 \ref{q:inner-ex5}}:例題 \ref{rei:inner-csapp} をまねて
$\vect{y} = (1, 2, 2)$ を選ぶ。
\item \textbf{問 \ref{q:inner-ex6}}:$\|\vect{a} \pm \vect{b}\|^2$ を
双線形性で展開する。(2) はコーシー・シュワルツの上限と比べる。
\item \textbf{問 \ref{q:inner-ex7}}:対称・双線形は成分計算、正値性は
$2a_1^2 + 3a_2^2 \ge 0$ から。係数が正であることが効く。
\item \textbf{問 \ref{q:inner-ex8}}:$\langle x^2 - c, 1 \rangle = 0$ を
$c$ について解く。
\item \textbf{問 \ref{q:inner-ex9}}:移項して
$\langle \vect{a}, \vect{b} - \vect{c} \rangle = 0$。$\vect{a}$ は
何でもよいのだから、いちばん役に立つものを選ぶ。
\item \textbf{問 \ref{q:inner-ex10}}:$\cos\theta_n = 1/\sqrt{n}$ になる。
$n \to \infty$ での挙動を考える。
\item \textbf{問 \ref{q:inner-ex11}}:分解して三角不等式を使うと
$\|\vect{a}\| - \|\vect{b}\| \le \|\vect{a} - \vect{b}\|$。
$\vect{a}$ と $\vect{b}$ を入れ替えてもう一本作る。
\item \textbf{問 \ref{q:inner-ex12}}:$\|\vect{a} + \vect{b}\|^2$ の展開式と
仮定を見比べる。
\end{itemize}
}

% =====================================================
\chapter{正規直交基底とグラム・シュミットの直交化法}\label{chap:gs}
% =====================================================

\section{正規直交基底のありがたみ}

\begin{teigi}[正規直交基底]
基底 $\vect{u}_1, \dots, \vect{u}_n$ が、互いに直交し
($i \neq j$ で $\langle \vect{u}_i, \vect{u}_j \rangle = 0$)、
各ノルムが $1$($\|\vect{u}_i\| = 1$)であるとき、
\textbf{正規直交基底}という。$\R^n$ の標準基底はその代表例である。
\end{teigi}

一般の基底では、ベクトル $\vect{v}$ の座標を求めるのに
連立方程式を解く必要があった。正規直交基底なら、
$\vect{v} = c_1\vect{u}_1 + \cdots + c_n\vect{u}_n$ の両辺と
$\vect{u}_i$ の内積をとるだけで、直交性により他の項が消え、
\[
c_i = \langle \vect{v}, \vect{u}_i \rangle
\]
と\textbf{座標が内積一発で求まる}。計算上も理論上も、
基底は正規直交であるに越したことはない。

\section{グラム・シュミットの直交化法}

幸い、どんな基底からでも正規直交基底を機械的に作れる。

\begin{teiri}[グラム・シュミットの直交化法]
線形独立な $\vect{a}_1, \dots, \vect{a}_k$ から、次の手順で
正規直交系 $\vect{u}_1, \dots, \vect{u}_k$ が得られる
(しかも各段階で張る空間は変わらない)。
\begin{enumerate}
\item $\vect{b}_1 = \vect{a}_1$ とする。
\item $\vect{b}_j = \vect{a}_j - \displaystyle\sum_{i=1}^{j-1}
\frac{\langle \vect{a}_j, \vect{b}_i \rangle}{\langle \vect{b}_i, \vect{b}_i \rangle}\,
\vect{b}_i$
(すでにできた方向への\textbf{影を引き去る})。
\item 最後に正規化:$\vect{u}_j = \vect{b}_j / \|\vect{b}_j\|$。
\end{enumerate}
\end{teiri}

手順 2 で引いている
$\dfrac{\langle \vect{a}, \vect{b} \rangle}{\langle \vect{b}, \vect{b} \rangle}\vect{b}$
は、$\vect{a}$ の $\vect{b}$ 方向への\textbf{正射影}
(高校で学んだ「正射影ベクトル」の式そのもの)である。
「新入りのベクトルから、既存メンバー方向の成分をすべて差し引けば、
残りは全員と直交する」という、素朴だが完璧なアイデアである。

\begin{rei}
$\vect{a}_1 = \mat{1 \\ 1 \\ 0}$, $\vect{a}_2 = \mat{0 \\ 1 \\ 1}$ を直交化する。
$\vect{b}_1 = \vect{a}_1$。
$\langle \vect{a}_2, \vect{b}_1 \rangle = 1$,
$\langle \vect{b}_1, \vect{b}_1 \rangle = 2$ より
\[
\vect{b}_2 = \mat{0 \\ 1 \\ 1} - \frac{1}{2}\mat{1 \\ 1 \\ 0}
= \mat{-1/2 \\ 1/2 \\ 1}.
\]
正規化して
\[
\vect{u}_1 = \frac{1}{\sqrt{2}}\mat{1 \\ 1 \\ 0}, \qquad
\vect{u}_2 = \frac{1}{\sqrt{6}}\mat{-1 \\ 1 \\ 2}.
\]
$\langle \vect{u}_1, \vect{u}_2 \rangle = 0$ を確かめてほしい。
\end{rei}

\section{直交補空間と正射影}

部分空間 $W$ に対し、$W$ のすべてのベクトルと直交するベクトルの全体
\[
W^{\perp} = \{\, \vect{v} \mid \text{すべての } \vect{w} \in W \text{ に対し }
\langle \vect{v}, \vect{w} \rangle = 0 \,\}
\]
を $W$ の\textbf{直交補空間}という。任意の $\vect{v}$ は
\[
\vect{v} = \underbrace{\vect{v}_W}_{\in W} +
\underbrace{\vect{v}_{\perp}}_{\in W^{\perp}}
\]
の形に\textbf{ただ一通り}に分解され、$\vect{v}_W$ を $\vect{v}$ の
$W$ への\textbf{正射影}という。

\begin{itembox}[l]{深掘り:正射影は「最良の近似」である}
正射影 $\vect{v}_W$ には、計算を超えた意味がある:
\textbf{$W$ の中で $\vect{v}$ に最も近い点}が $\vect{v}_W$ なのである。
実際、任意の $\vect{w} \in W$ に対し、三平方の定理により
\[
\|\vect{v} - \vect{w}\|^2 = \|\vect{v} - \vect{v}_W\|^2 + \|\vect{v}_W - \vect{w}\|^2
\ge \|\vect{v} - \vect{v}_W\|^2
\]
となる($\vect{v} - \vect{v}_W \in W^{\perp}$ と
$\vect{v}_W - \vect{w} \in W$ が直交することを使った)。

「届かない点には、垂線の足で最善を尽くす」――この原理が、
データ分析の根幹をなす\textbf{最小二乗法}の正体である。
観測データに厳密に合う解が存在しないとき
($A\vect{x} = \vect{b}$ が解なしのとき)、
$\vect{b}$ を $A$ の列が張る空間へ正射影した点なら必ず到達でき、
それが「誤差最小の近似解」になる。第VII部で全面展開する。
\end{itembox}

\begin{mondai}\label{q:gs}
$\vect{a}_1 = \mat{1 \\ 0 \\ 1}$, $\vect{a}_2 = \mat{1 \\ 1 \\ 0}$ に
グラム・シュミットの直交化法を適用し、正規直交系を作れ。
\end{mondai}

% =====================================================
\chapter{直交行列}\label{chap:orth}
% =====================================================

\begin{teigi}[直交行列]
${}^t\!P\,P = E$、すなわち ${}^t\!P = P^{-1}$ を満たす正方行列 $P$ を
\textbf{直交行列}という。
\end{teigi}

定義の式 ${}^t\!P P = E$ を成分で読むと、
「$P$ の列ベクトルたちが正規直交系をなす」ことと同値である。
つまり\textbf{直交行列とは、正規直交基底を列に並べた行列}である
(逆行列が転置で求まるのはそのご褒美である)。

\begin{teiri}[直交行列の性質]
$P$ を直交行列とすると、
\begin{enumerate}
\item 内積を保つ:$\langle P\vect{x}, P\vect{y} \rangle = \langle \vect{x}, \vect{y} \rangle$。
したがって長さと角度を保つ。
\item $\det P = \pm 1$。
\item 直交行列の積、逆行列も直交行列である。
\end{enumerate}
\end{teiri}

性質 1 は $\langle P\vect{x}, P\vect{y} \rangle = {}^t\vect{x}\,{}^t\!P P\,\vect{y}
= {}^t\vect{x}\,\vect{y}$ から、性質 2 は
${}^t\!PP = E$ の行列式をとり $(\det P)^2 = 1$ から従う。

長さと角度を保つ変換とは、図形を\textbf{変形させない}変換である。
平面($2$ 次)の直交行列は、回転($\det = 1$)と鏡映($\det = -1$)で
すべて尽くされることが知られている。第I部に登場した回転行列 $R_\theta$ が
直交行列であることは、列の正規直交性からすぐ確かめられる。

\begin{mondai}\label{q:orthm}
　(1) 回転行列 $R_\theta$ が直交行列であることを ${}^t R_\theta R_\theta = E$ の
計算で確かめよ。\\
　(2) 直交行列 $P, Q$ の積 $PQ$ が直交行列であることを示せ。
\end{mondai}

% =====================================================
\chapter{対称行列のスペクトル分解}\label{chap:spectral}
% =====================================================

\section{対称行列の固有値・固有ベクトルの奇跡}

${}^t\!A = A$ を満たす行列を\textbf{対称行列}という。
対称行列の固有値理論は、一般の行列では望めない完璧さをもつ。

\begin{teiri}[対称行列の基本性質]
実対称行列 $A$ について、
\begin{enumerate}
\item 固有値はすべて\textbf{実数}である。
\item 相異なる固有値に属する固有ベクトルは\textbf{直交}する。
\end{enumerate}
\end{teiri}

性質 2 の証明は美しい。$A\vect{u} = \lambda\vect{u}$,
$A\vect{v} = \mu\vect{v}$($\lambda \neq \mu$)とすると、
対称性 ${}^t\!A = A$ により
\[
\lambda \langle \vect{u}, \vect{v} \rangle
= \langle A\vect{u}, \vect{v} \rangle
= {}^t\vect{u}\,{}^t\!A\,\vect{v}
= \langle \vect{u}, A\vect{v} \rangle
= \mu \langle \vect{u}, \vect{v} \rangle
\]
となり、$\lambda \neq \mu$ から $\langle \vect{u}, \vect{v} \rangle = 0$。
\textbf{内積を行列の左右に渡せる}ことが対称行列の本質である。

\begin{teiri}[実対称行列のスペクトル定理]
実対称行列 $A$ は、\textbf{直交行列} $P$ によって対角化できる:
\[
{}^t\!P A P = \mat{\lambda_1 & & \\ & \ddots & \\ & & \lambda_n}.
\]
すなわち、固有ベクトルからなる\textbf{正規直交基底}が必ず存在する。
固有値が重複していても、対角化に失敗することはない。
\end{teiri}

第IV部では「対角化できない行列がある」ことに悩まされたが、
対称行列の世界ではその心配が消える。しかも基底変換は直交行列、
つまり\textbf{回転(と鏡映)だけ}で済む。
変換 $A$ は「直交する $n$ 本の軸に沿って、それぞれ $\lambda_i$ 倍に
伸縮するだけ」という、考えうる最も素直な姿に分解される。

\begin{rei}
$A = \mat{2 & 1 \\ 1 & 2}$。固有値は $3, 1$、固有ベクトルは
$\mat{1 \\ 1}$, $\mat{1 \\ -1}$(確かに直交している)。正規化して並べた
\[
P = \frac{1}{\sqrt{2}}\mat{1 & 1 \\ 1 & -1}
\]
は直交行列で、${}^t\!PAP = \mat{3 & 0 \\ 0 & 1}$ となる。
\end{rei}

スペクトル定理は和の形にも書ける。固有値 $\lambda_i$ の固有空間への
正射影行列を $P_i$ とすると
\[
A = \lambda_1 P_1 + \lambda_2 P_2 + \cdots + \lambda_k P_k
\]
(\textbf{スペクトル分解})。「行列 = 固有値 × 射影 の重ね合わせ」という
この見方は、量子力学(観測値 = 固有値、観測 = 射影)の数学的骨格そのものである。

\begin{itembox}[l]{深掘り:なぜ対称行列はこれほど頻繁に現れるのか}
対称行列は、応用数学の主要な舞台ほぼすべてに現れる。
(i) 多変数関数の $2$ 階偏導関数を並べた\textbf{ヘッセ行列}は、
偏微分の順序交換($f_{xy} = f_{yx}$)により対称になる。
(ii) 統計学の\textbf{分散共分散行列}は、共分散の定義の対称性
($\mathrm{Cov}(X, Y) = \mathrm{Cov}(Y, X)$)により対称になる
(その固有ベクトルを取り出す手法が主成分分析である。第IX部で扱う)。
(iii) 力学の慣性テンソル、量子力学の物理量(エルミート行列は対称行列の複素版)。
共通する理由は、これらが本質的に\textbf{二つの方向の対をとる操作}
(微分を $2$ 回、変数の組の共分散)から生まれ、対の順序に意味がないことである。
「世界の二次的な構造は対称行列で書かれる」と言ってよい。
だからこそ、対称行列でだけ成り立つ完璧な定理(スペクトル定理)の
応用範囲は絶大なのである。
\end{itembox}

\begin{mondai}\label{q:spec}
$A = \mat{2 & 1 \\ 1 & 2}$ について、固有空間への正射影行列
$P_1, P_2$ を求め、スペクトル分解 $A = 3P_1 + 1 \cdot P_2$ が
成り立つことを確かめよ。
(ヒント:単位固有ベクトル $\vect{u}$ への正射影行列は $\vect{u}\,{}^t\vect{u}$。)
\end{mondai}

% =====================================================
\chapter{二次形式}\label{chap:quadform}
% =====================================================

\section{二次形式と対称行列}

\begin{teigi}[二次形式]
$n$ 変数の $2$ 次の同次多項式を\textbf{二次形式}という。
たとえば $2$ 変数では $q(x, y) = ax^2 + 2bxy + cy^2$。
これは対称行列を使って
\[
q(x, y) = \mat{x & y}\mat{a & b \\ b & c}\mat{x \\ y}
= {}^t\vect{x} A \vect{x}
\]
と書ける。交差項 $2bxy$ の係数を半分ずつ $(1,2)$ 成分と $(2,1)$ 成分に
振り分けるのが要点である。
\end{teigi}

\section{標準形 ― 交差項を消す}

二次形式の解析を阻むのは交差項 $xy$ である。
スペクトル定理がこれを一掃する。$A$ を直交行列 $P$ で対角化し、
変数変換 $\vect{x} = P\vect{y}$ を行うと
\[
{}^t\vect{x}A\vect{x} = {}^t\vect{y}\,({}^t\!PAP)\,\vect{y}
= \lambda_1 y_1^2 + \lambda_2 y_2^2 + \cdots + \lambda_n y_n^2.
\]
\textbf{交差項のない標準形}が得られた。
しかも $P$ は直交行列(座標軸の回転)なので、図形の形は歪まない。

\begin{rei}[二次曲線の分類]
曲線 $2x^2 + 2xy + 2y^2 = 1$ は何か。行列 $\mat{2 & 1 \\ 1 & 2}$ の
固有値は $3, 1$ だから、軸を $45^\circ$ 回転した座標 $(X, Y)$ で
\[
3X^2 + Y^2 = 1
\]
となる。これは\textbf{楕円}である(長軸・短軸の方向が固有ベクトルの方向、
半径が $1/\sqrt{\lambda_i}$)。高校で「$xy$ 項のある $2$ 次曲線」が
扱えなかった理由と、その完全な解決がここにある。
固有値の符号が $(+, +)$ なら楕円、$(+, -)$ なら双曲線
――\textbf{曲線の種類は固有値の符号だけで決まる}。
\end{rei}

\section{正定値性}

\begin{teigi}[正定値]
二次形式 $q(\vect{x}) = {}^t\vect{x}A\vect{x}$($A$ は対称)が、
$\vect{x} \neq \vect{0}$ のすべてで $q(\vect{x}) > 0$ となるとき、
$q$(および $A$)は\textbf{正定値}であるという。
すべてで $q(\vect{x}) < 0$ なら\textbf{負定値}、
正負両方の値をとるなら\textbf{不定値}という。
\end{teigi}

標準形を見れば判定は明快である。

\begin{teiri}[正定値性の判定]
対称行列 $A$ について、次は同値である。
\begin{enumerate}
\item $A$ は正定値である。
\item $A$ の固有値はすべて正である。
\item $A$ の首座小行列式(左上の $k \times k$ 部分の行列式、
$k = 1, \dots, n$)がすべて正である(\textbf{シルベスターの判定法})。
\end{enumerate}
\end{teiri}

条件 3 は固有値を計算せずに済む実用的な判定法である。
$2$ 次なら「$a > 0$ かつ $\det A > 0$」だけでよい。

\begin{itembox}[l]{深掘り:極値判定の正体 ― 一変数から多変数へ}
一変数関数の極値判定「$f'(p) = 0$ かつ $f''(p) > 0$ なら極小」は
高校以来の常識である。多変数ではどうなるか。
$2$ 変数関数 $f(x, y)$ を停留点($f_x = f_y = 0$ となる点)のまわりで
テイラー展開すると、$1$ 次の項が消えるため、増分はおよそ
\[
f(p + \vect{h}) - f(p) \approx \frac{1}{2}\,{}^t\vect{h}\,H\,\vect{h},
\qquad
H = \mat{f_{xx} & f_{xy} \\ f_{xy} & f_{yy}}
\]
という\textbf{二次形式}で支配される($H$ はヘッセ行列)。よって
\[
H \text{ が正定値} \Rightarrow \text{極小}, \qquad
H \text{ が負定値} \Rightarrow \text{極大}, \qquad
H \text{ が不定値} \Rightarrow \text{鞍点}.
\]
微積分の教科書に載っている判定条件
「$f_{xx} > 0$ かつ $f_{xx}f_{yy} - f_{xy}^2 > 0$ なら極小」は、
シルベスターの判定法そのものである。

この話は経済学で決定的に重要になる。効用関数や生産関数の\textbf{凹性}
(ヘッセ行列が至るところ負定値ないし半負定値)は、
最適化問題の解が真に最大値であることを保証する条件であり、
ミクロ経済学の理論全体がこの上に建っている。
最適化・機械学習でも、損失関数のヘッセ行列の固有値分布が
学習の難しさを左右する。\textbf{「最適化の言語は二次形式である」}。
\end{itembox}

\begin{mondai}\label{q:quad}
　(1) 二次形式 $q(x, y) = 2x^2 + 2xy + 2y^2$ が正定値であることを、
シルベスターの判定法で確かめよ。\\
　(2) 曲線 $x^2 + 4xy + y^2 = 1$ は楕円か双曲線か。
固有値を求めて判定し、標準形を書け。
\end{mondai}
